%! Competing Function Model Validation — Unit 4, Lesson 4.7 — Homework (Answer Key)
\documentclass[10pt]{article}
\usepackage{precalculus-article}
\usepackage{precalculus-key}   % pulls in -boxes; do NOT also load -boxes
\newcommand{\TallMath}[1]{$\displaystyle #1\rule[-1.4em]{0pt}{3.2em}$}

\begin{document}

\pageheader{Unit 4, Lesson 4.7}{Homework --- Answer Key}
\namedateperiod

\vspace{0.1in}

\begin{teachernote}
Problem 1: $\Delta_2$ values are constant at $2$, so quadratic. Note: the ratios are NOT
constant ($5/2=2.5$, $10/5=2$, $17/10=1.7$, $26/17\approx1.53$), confirming not exponential.
Problem 2: Linear $\hat{y}=79.8$, Quadratic $\hat{y}=78.6$, Exponential $\hat{y}=4\cdot1.7^6=4\cdot24.14\approx96.6$.
Errors: Linear $+2.2$, Quadratic $+3.4$, Exponential $-14.6$. Smallest: Linear.
Problem 5: Answer is (C). Problem 6: SSR(Linear) = 9+1+4+64+225 = 303; SSR(Exp) = 0.25+0.04+0.01+0.09+0.16 = 0.55.
\end{teachernote}

\begin{enumerate}[itemsep=16pt]

\item
      \renewcommand{\arraystretch}{1.7}
      \begin{tabular}{c|c|c|c|c}
      $x$ & $y$ & $\Delta_1$ & $\Delta_2$ & Ratio \\ \hline
      0 &  2 & \ans{3} & \ans{2} & \ans{2.5} \\
      1 &  5 & \ans{5} & \ans{2} & \ans{2.0} \\
      2 & 10 & \ans{7} & \ans{2} & \ans{1.7} \\
      3 & 17 & \ans{9} & --- & \ans{1.53} \\
      4 & 26 & --- & --- & --- \\
      \end{tabular}

      \vspace{0.04in}
      The data is best modeled by a(n) \ans{quadratic} function. Explain: \ans{The second differences are constant ($\Delta_2 = 2$), which is the signature of a quadratic model.}

\item
      \begin{center}
      \renewcommand{\arraystretch}{1.4}
      \begin{tabular}{l|c|c|c}
      \textbf{Model} & \textbf{Equation} & \textbf{Predicted $\hat{y}$ at $x=6$} & \textbf{Error} \\ \hline
      Linear      & $\hat{y} = 13.5x - 1.2$       & \ans{79.8}  & \ans{$+2.2$} \\
      Quadratic   & $\hat{y} = 2.1x^2 + 0.5x - 3$ & \ans{78.6}  & \ans{$+3.4$} \\
      Exponential & $\hat{y} = 4 \cdot 1.7^x$     & \ans{$\approx 96.6$} & \ans{$\approx -14.6$} \\
      \end{tabular}
      \end{center}

      Model with smallest absolute error at $x = 6$: \ans{Linear ($|{+2.2}| = 2.2$)}

\item
      \begin{enumerate}[label=(\alph*), itemsep=6pt]
        \item Most appropriate model: \ans{Exponential}
        \item \ansline{Doubling every hour means the output is multiplied by 2 for each unit increase in the input --- a constant successive ratio of 2 --- which is the defining characteristic of an exponential model.}
      \end{enumerate}

\item
      \begin{enumerate}[label=(\alph*), itemsep=6pt]
        \item Which plot suggests the model is appropriate? \ans{Plot B}
        \item \ansline{Plot A's U-shaped pattern indicates the chosen model is the wrong function type; the model is systematically under- and over-predicting in a structured way, suggesting a different function family should be used.}
      \end{enumerate}

\item Answer: \ans{(C)} --- Exponential has both the highest $R^2$ (0.99) and a patternless residual plot, satisfying both fit criteria simultaneously.

\item
      \begin{enumerate}[label=(\alph*), itemsep=6pt]
        \item SSR (Linear): $3^2 + 1^2 + (-2)^2 + (-8)^2 + (-15)^2 = 9 + 1 + 4 + 64 + 225 = $ \ans{303}

              SSR (Exponential): $0.5^2 + (-0.2)^2 + 0.1^2 + (-0.3)^2 + 0.4^2 = 0.25 + 0.04 + 0.01 + 0.09 + 0.16 = $ \ans{0.55}

        \item \ansline{The exponential model has a much smaller SSR (0.55 vs.\ 303). A smaller SSR means the model's predictions are closer to the actual values overall, indicating a much better fit.}
      \end{enumerate}

\end{enumerate}

\end{document}
